%% kapa: exam=1 round=32 year=2021 period=1 form=A subject=economics subject_ko=경제학원론 seq=2 number=42
%% source: Q-Net 감정평가사 기출문제 (cst00309) artlSeq=5208397 file=extracted/1008/5208397/A형/2021년도 제32회 감정평가사 1차 1교시 A형.pdf
%% conversion: kordoc 4.24.51 -> tools/gen1.py (auto); review=proofread-sonnet; flags=none
\begin{question}{42}
사회후생에 관한 설명으로 옳지 않은 것은?
\begin{choices}
  \item 차선의 이론은 부분적 해결책이 최적은 아닐 수 있음을 보여준다.
  \item 롤즈(J. Rawls)적 가치판단을 반영한 사회무차별곡선은 L자 모양이다.
  \item 파레토 효율성 조건은 완전경쟁의 상황에서 충족된다.
  \item 공리주의적 사회후생함수는 최대다수의 최대행복을 나타낸다.
  \item 애로우(K. Arrow)의 불가능성 정리에서 파레토원칙은 과반수제를 의미한다.
\end{choices}
\end{question}
\answer{5}
